\section{Implementation Status and Open Problems}
\label{sec:status}

This section describes the reference implementation, reports the status of each major mechanism, and lists unresolved architectural questions.

\subsection{Implementation}

The Freenet Core is implemented in Rust as an asynchronous \texttt{tokio}-based runtime. Contracts and delegates execute inside Wasmtime \citep{wasmtime}; both run under per-instance fuel limits and explicit memory bounds so that a runaway or malicious module cannot exhaust the host. Contract state is persisted through a pluggable store layer with implementations on top of \texttt{redb} \citep{redb} and SQLite.

The runtime is organized as a single central event loop, a \texttt{tokio::select!} over inbound transport packets, client API events, operation timers, and background tasks, that dispatches into per-operation state machines. The single-coordinator design simplifies reasoning about partial state but limits the throughput of a single peer to what one event-loop thread can sustain. So far this has not been the binding constraint in deployment.

\paragraph{Deterministic simulation testing.} A peer-to-peer system is hard to test by ordinary means: many bugs in routing, congestion, or convergence appear only in specific multi-peer interleavings that are difficult to reproduce. The Core ships with a deterministic simulation harness built on Turmoil \citep{turmoil}: virtual time, a seeded global RNG, an in-memory socket implementation, and fault-injection knobs for loss, latency, partitions, and reorderings. Tests run the entire network in a single deterministic process and replay bit-for-bit from a seed. Most non-trivial routing, subscription, and contract-execution bugs we have found were first reproduced in simulation.

\paragraph{Content-addressed components.} Contracts, delegates, and UI bundles are all addressed by hash. This pushes a lot of complexity off the platform: there is no version negotiation, no upgrade protocol, no need for the platform to understand application semantics during deployment. A new version is published under a new key; the old key continues to work as long as anyone hosts it.

\subsection{Implementation Status of Each Mechanism}

The following table reports the current status of the mechanisms described in this paper. ``Deployed'' means present in the production runtime and exercised on the live network; ``Implemented'' means present in the runtime but not yet broadly exercised at scale; ``Experimental'' means available as a build-time option but not in production; ``Open'' means designed or partially designed but not implemented.

\begin{center}
\small
\begin{tabular}{p{0.42\textwidth}p{0.18\textwidth}p{0.32\textwidth}}
\toprule
Mechanism & Status & Notes \\
\midrule
Small-world ring topology & Deployed & Live network $\approx 440$ active peers \\
Greedy forwarding, HTL-bounded & Deployed & HTL cap = 10; greedy/performance-ranked every hop \\
Two-stage CONNECT acceptance & Deployed & Greedy forward, accept-and-forward in 5\% band, Kleinberg-scored at terminus \\
Adaptive isotonic-regression routing & Deployed & Per-neighbor distance-indexed model \\
Summary/delta synchronization & Deployed & Per-contract, application-defined \\
Subscription trees with leases & Deployed & 8-min lease, 2-min renewal, interest-gated \\
Neighbor host-advertisement mesh & Deployed & Update fan-out and terminus lookup \\
Contract WebAssembly execution & Deployed & Wasmtime, fuel + memory limits \\
Delegates (local, secret state) & Deployed & Used by River chat \citep{river} \\
Delegate cross-device sync via shared-secret contract & Partial & Encrypted-contract primitives exist; end-to-end delegate flow not wired up \\
NAT hole-punching transport & Deployed & UDP, X25519 + AEAD \\
Fixed-rate congestion control & Deployed & Production default \\
LEDBAT++ congestion control & Experimental & Did not meet our requirements in initial trials; not in production \\
BBR-style congestion control & Experimental & Build-time option, not the default \\
Deterministic simulation testing & Deployed & Turmoil-based \citep{turmoil} \\
Per-peer resource limits (bandwidth, storage) & Deployed & Capability-relative default budget, operator-overridable \\
Sybil-resistant peer identity / location & Open & Defenses are application-layer (see \cref{sec:trust}) \\
Network-level incentive / abuse layer & Open & Application-layer reputation contracts are the intended mechanism \\
\bottomrule
\end{tabular}
\end{center}

The single biggest gap between this paper's narrative and the deployed system is congestion control. The platform supports pluggable congestion controllers but the deployed default is a token-bucket fixed-rate scheme. LEDBAT++ and BBR variants exist in the codebase as alternatives, but neither has yet been adapted to the workload satisfactorily. Real congestion control across heterogeneous peer-to-peer links remains open work.

\subsection{Where the Platform Fits}

Contract-defined merge, summary/delta synchronization, small-world routing, and delegates support applications whose conflicts can be resolved by application-defined merge. Examples include group chat, collaborative documents, social feeds, reputation systems, asynchronous inboxes, and identity directories. A working chat implementation \citep{river} and an anonymous identity scheme \citep{ghostkey} are deployed on the live network.

\subsection{Where the Platform Does Not Fit}

Several application classes remain poorly suited to the model.

\paragraph{Strictly linearizable transactions.} A defense against double-spending fungible tokens requires that two attempts to spend the same balance cannot both succeed. The merge model cannot express ``one of these updates is rejected unconditionally''; merge is total. Concretely, fungible-token transfer with platform-level double-spend protection cannot be expressed by a single contract in the current model. Applications that need a strong-consistency primitive must either contain the inconsistency inside a contract whose validity predicate detects and rejects the second update (which forces the conflict into application code, with the loser's state restored on detection) or layer an auxiliary consensus mechanism on top of the platform.

The limitation extends beyond currency. A reputation system that only \emph{accumulates} endorsements is monotone and fits the merge model. Revoking a credential, slashing a stake, or reducing a balance introduces the same requirement to reject a conflicting second update. Applications that otherwise fit the model may therefore require an additional consistency mechanism for non-monotone operations.

\paragraph{Abuse-resistant open posting.} A purely-open contract is simple to write as a set under union and equally simple to flood. Practical contracts in this class restrict writes to signed records by approved keys, charge a proof-of-work fee per update, or require credentials issued by a reputation contract. The platform's transport-layer rate limiting and per-neighbor resource accounting do provide a partial floor: a neighbor consuming excessive resources is disconnected, which caps the volume any single attacker can push through any single peer. That bound is useful but limited; it does not stop a distributed flood from many cheap identities, nor a slow-burn pattern that stays within per-peer limits. The platform has no application-level spam defense beyond that transport floor.

\paragraph{Coordinated DoS.} Per-peer resource limits protect a peer against a single misbehaving neighbor, not against coordinated traffic from many neighbors using cheap identities. Defenses here probably require accounting that persists across identities, which is the open incentive layer noted above.

\paragraph{Durable storage of cold state.} Contract state persists only while at least one peer near $\ell(\contract)$ retains it. Each peer holds contracts in a byte-budgeted hosted set whose size defaults to a fraction of host memory. When over budget, the peer evicts the contract with the fewest dependent subscribers, prioritizing local client subscriptions over downstream subscribers and breaking ties by least-recent genuine \textsf{GET} or \textsf{PUT} access. Subscriber count affects eviction order but does not pin a contract. If necessary, the peer evicts a subscribed contract and removes its subscription.

Eviction removes the persisted state, parameters, and code blob from disk. Cold contracts with no readers or subscribers are therefore the first candidates. The platform does not recruit a replacement when a host evicts a contract; recovery requires a new \textsf{PUT} from a party that still holds the state. A contract may consequently disappear from every host and become unrecoverable.

The platform provides no durability guarantee for cold state. Dormant collaborative documents, unpolled inboxes, and identity records for inactive users may remain cold for long periods. Maintaining a minimum replica count independently of demand---through active repair, incentivized pinning, or archival contracts---remains open work.

\subsection{Open Problems}

The following questions require further implementation or empirical evaluation.

\begin{enumerate}[leftmargin=1.4em, itemsep=2pt]
  \item \textbf{Scaling characterization.} The single 24-hour data point in \cref{sec:routing-measurement} is consistent with logarithmic-diameter routing at the deployed network size. Characterization requires measurements across orders of magnitude in $N$. The current deployment cannot provide that range, and the deterministic simulation harness must be extended to emit hop-count-versus-$N$ data before it can serve as an alternative.
  \item \textbf{Failure modes of contract-defined merge.} A contract whose merge is not associative or idempotent does not converge. The platform exposes the resulting divergence but takes no corrective action. Better diagnostics and, where feasible, platform-checked contract invariants remain open work. A near-term improvement would be a property-based test harness that generates states and updates to exercise associativity, commutativity, idempotence, and sync soundness (\cref{prop:sync}). Such a harness cannot constrain a malicious author, but it would provide honest authors with a repeatable CI workflow.
  \item \textbf{Congestion control.} The fixed-rate default remains temporary while the alternative controllers are adapted and evaluated for this workload.
  \item \textbf{Sybil-resistant identity.} Ghost Keys demonstrates an application-layer approach, but no general mechanism exists. Identity and reputation contracts are routed through the location layer they are intended to secure. An adversary with enough address prefixes to crowd a contract's location can serve poisoned reads and prevent honest writes from reaching its canonical replicas (\cref{sec:trust}). Prefix collisions also reduce the geographic diversity of honest peers. Resolving this circular dependency requires ring-distance separation among replicas of safety-critical contracts, demonstrable geographic diversity near their locations, or out-of-band trust anchors for the contracts that bootstrap the rest.
  \item \textbf{Incentive / abuse layer.} The network needs a mechanism to resist coordinated denial of service without centralized admission control or excessive barriers to joining.
  \item \textbf{Durability of cold state.} A contract that no one reads, subscribes to, or updates can be evicted from every host. A minimum replica count maintained through active repair, incentivized pinning, or archival contracts must avoid creating an unbounded free-storage abuse surface.
\end{enumerate}
